\chapter{Protein Electrophoresis and Immunofixation}

\section*{What Is This Test?}
Serum protein electrophoresis (SPEP) is a laboratory technique that separates serum proteins into distinct fractions based on their electrophoretic mobility (charge and size): albumin, alpha-1 globulins, alpha-2 globulins, beta globulins, and gamma globulins. Immunofixation electrophoresis (IFE) is a complementary technique performed on the same specimen that identifies the specific immunoglobulin heavy chain class (IgG, IgA, IgM, IgD, IgE) and light chain type (kappa or lambda) of monoclonal proteins detected by SPEP. Together, SPEP and IFE form the cornerstone diagnostic tests for monoclonal gammopathies. A monoclonal gammopathy occurs when a single clone of plasma cells (or B cells) proliferates and produces a homogenous (monoclonal) immunoglobulin or immunoglobulin fragment. This monoclonal protein (M-protein, M-spike, or paraprotein) appears as a sharp, discrete band on SPEP, typically in the gamma region (sometimes in the beta or alpha-2 region for IgA or light chain disease). IFE confirms clonality by demonstrating that the M-protein is composed of a single immunoglobulin heavy chain class and a single light chain type (e.g., IgG kappa, IgA lambda). SPEP and IFE are required for the diagnosis and monitoring of monoclonal gammopathy of undetermined significance (MGUS), multiple myeloma, Waldenstrom macroglobulinemia, AL amyloidosis, solitary plasmacytoma, POEMS syndrome, and other plasma cell and lymphoproliferative disorders.

\section*{Alternative Names}
SPEP; Serum Protein Electrophoresis; IFE; Immunofixation Electrophoresis; Immunofixation; Protein Electrophoresis; M-Protein Screen; Monoclonal Protein Study; Paraprotein Screen; SPE; Immunotyping

\section*{Principle}
Serum protein electrophoresis is performed on agarose gel or capillary electrophoresis systems. Serum is applied to a gel matrix, and an electric current is applied. Serum proteins migrate through the gel at different rates depending on their net charge and size. After a fixed time, the proteins are fixed, stained (typically with Coomassie brilliant blue or amido black), destained, and scanned by a densitometer. The densitometer trace produced shows five major peaks (fractions): albumin (the largest peak, 52--65\% of total protein), alpha-1 globulins (alpha-1-antitrypsin, alpha-1-acid glycoprotein, 2--5\%), alpha-2 globulins (alpha-2-macroglobulin, haptoglobin, ceruloplasmin, 6--12\%), beta globulins (transferrin, beta-lipoprotein/LDL, C3, 7--15\%), and gamma globulins (immunoglobulins, 10--20\%). The M-protein appears as a sharp, narrow peak (spike) anywhere from the late alpha-2 to the gamma region. The area under the M-spike is calculated by the densitometer to quantify the M-protein concentration. For immunofixation electrophoresis (IFE), serum proteins are first separated by agarose gel electrophoresis. Then, specific antisera directed against each immunoglobulin heavy chain (anti-gamma, anti-alpha, anti-mu, anti-delta, anti-epsilon) and each light chain type (anti-kappa, anti-lambda) are applied to separate lanes of the gel. If a monoclonal protein is present, it will form a discrete precipitin band with only one heavy chain antiserum and only one light chain antiserum (e.g., a band in the IgG lane and the kappa lane at the same electrophoretic position), confirming clonality. A polyclonal increase shows diffuse staining of both kappa and lambda lanes. Capillary zone electrophoresis (CZE) is an alternative method that separates proteins in free solution in a capillary tube under high voltage, with UV detection, providing a similar electrophoretogram \cite{dispenzieri2009standardization}.

\section*{History}
Serum protein electrophoresis was developed by Arne Tiselius in 1937, for which he received the Nobel Prize in Chemistry in 1948. Tiselius used a moving boundary electrophoresis apparatus (the Tiselius apparatus) to separate serum proteins into albumin, alpha, beta, and gamma globulins. Zone electrophoresis on paper and later cellulose acetate, starch gel, and agarose gel was developed in the 1950s--1960s. Immunofixation electrophoresis was developed by Alper and Johnson in 1969 and refined for clinical use by Cawley and Ritchie. The introduction of automated agarose gel electrophoresis and capillary electrophoresis systems in the 1980s--1990s greatly improved throughput and precision. In 2014, the International Myeloma Working Group updated the diagnostic criteria for multiple myeloma and endorsed the combination of SPEP, IFE, and serum free light chain assay as the initial screening panel for monoclonal gammopathies \cite{rajkumar2014international}. In 2015, Katzmann et al. demonstrated that replacing 24-hour urine studies with serum IFE and FLC for initial monoclonal gammopathy screening had a negative predictive value >99\%, simplifying the workup for most patients \cite{katzmann2009screening}.

\section*{Why This Test Is Ordered}
\begin{itemize}
  \item Screening for monoclonal gammopathy in patients with unexplained anemia, hypercalcemia, renal insufficiency, bone pain, pathologic fractures, peripheral neuropathy, recurrent infections, hyperviscosity symptoms, or unexplained proteinuria
  \item Diagnosis and classification of monoclonal gammopathies: monoclonal gammopathy of undetermined significance (MGUS, M-protein <3 g/dL, <10\% bone marrow plasma cells, no CRAB criteria), multiple myeloma (M-protein $\geq$3 g/dL in serum or $\geq$500 mg/24 hours in urine, $\geq$10\% clonal plasma cells in bone marrow, with or without CRAB criteria or myeloma-defining events), smoldering multiple myeloma (M-protein $\geq$3 g/dL and/or $\geq$10\% bone marrow plasma cells without CRAB criteria or myeloma-defining events)
  \item Diagnosis of Waldenstrom macroglobulinemia (IgM monoclonal gammopathy with bone marrow infiltration by lymphoplasmacytic lymphoma)
  \item Evaluation for AL amyloidosis in a patient with nephrotic-range proteinuria, unexplained heart failure with preserved ejection fraction, hepatomegaly, macroglossia, or unexplained peripheral neuropathy
  \item Monitoring of treatment response and detection of relapse in multiple myeloma and other monoclonal gammopathies. Measurement of M-protein by SPEP is the primary method for assessing response per IMWG uniform response criteria
  \item Evaluation of hypogammaglobulinemia (decreased gamma globulin fraction) as a marker of immunodeficiency
  \item Detection of polyclonal hypergammaglobulinemia (diffuse increase in gamma globulins) in chronic inflammatory diseases, autoimmune diseases, cirrhosis, and chronic infections
  \item Detection of acute phase protein changes: decreased albumin with elevated alpha-1 and alpha-2 globulins in acute inflammation
\end{itemize}

\section*{Primary vs Secondary Test}
This is a primary diagnostic test for plasma cell dyscrasias and monoclonal gammopathies. SPEP with IFE, together with the serum free light chain assay, is the recommended initial screening panel for all suspected monoclonal gammopathies per IMWG guidelines.

\section*{How This Test Is Performed}
For most patients, a health care professional will take a blood sample from a vein in your arm using a small needle. After the needle is inserted, a small amount of blood will be collected into a test tube or vial. You may feel a little sting when the needle goes in or out. This usually takes less than five minutes.

\section*{What Preparation Is Needed}
No special preparation is required. Fasting is not necessary. Certain medications that can produce a monoclonal or oligoclonal banding pattern should be disclosed: intravenous immunoglobulin (IVIG) therapy produces a transient polyclonal or oligoclonal banding pattern that can be mistaken for an M-protein on SPEP and IFE. Fibrinogen (present in plasma samples but not serum) can produce a monoclonal-like band in the beta-gamma region; this is one reason serum is preferred over plasma for SPEP/IFE. Patients receiving monoclonal antibody therapies (daratumumab, isatuximab, elotuzumab) produce a small IgG kappa M-spike on SPEP and IFE that represents the therapeutic monoclonal antibody, not endogenous disease \cite{mccudden2010interference}. This drug-related M-spike can mimic residual disease and confound response assessment. Laboratories performing SPEP on patients receiving these drugs should be alerted so that appropriate measures (e.g., daratumumab-specific IFE reflex testing or mass spectrometry-based monoclonal protein detection) can be performed.

\section*{Contraindications and Precautions}
No absolute contraindications. Specimen must be collected in a serum separator tube (SST, gold-top) or red-top tube. Plasma (EDTA, heparin, citrate) should NOT be used because fibrinogen produces a monoclonal-like band in the beta-gamma region that can be misinterpreted as an M-protein. Hemolyzed specimens interfere with electrophoretic separation and quantification of alpha and beta globulins. The sensitivity of SPEP for detecting low-concentration monoclonal proteins is approximately 0.05--0.1 g/dL (500--1,000 mg/L). M-proteins below this concentration (small M-proteins) may not be visible on SPEP but may be detected by IFE or free light chain assay. Monoclonal free light chains (Bence Jones proteinemia) are poorly detected by SPEP because free light chains are rapidly cleared by the kidneys; the serum free light chain assay is the preferred method for detecting light chain disease. Filter paper support media (for protein electrophoresis) must be free of protein contamination. Capillary electrophoresis instruments must be properly maintained to prevent protein carryover and baseline artifacts. Therapeutic monoclonal antibodies (daratumumab, isatuximab, elotuzumab) produce an IgG kappa band on IFE that can interfere with assessment of complete response in patients with IgG kappa myeloma. The laboratory must be notified, and the drug band must be accounted for using specialized methods.

\section*{Risks}
Minimal risk. Slight pain or bruising at the venipuncture site. Rare: hematoma, infection, vasovagal syncope.

\section*{What Happens After the Test}
Results are typically available within 1 to 3 days. The SPEP report includes the densitometer tracing showing protein fractions, the M-protein concentration (if present, measured as the area under the spike in g/dL), and the interpretation. IFE results confirm clonality and identify the immunoglobulin type. An abnormal SPEP with confirmed monoclonal protein on IFE should prompt further evaluation to determine whether the monoclonal gammopathy represents a benign condition (MGUS) or a malignant plasma cell disorder (multiple myeloma, Waldenstrom macroglobulinemia, AL amyloidosis).

\section*{Understanding Results}

\subsection*{Negative SPEP (No M-Protein Detected)}
\begin{itemize}
  \item No monoclonal protein (M-spike) on SPEP. Approximately 15--20\% of patients with multiple myeloma (particularly light chain myeloma or non-secretory myeloma) have a negative SPEP. The serum IFE and serum free light chain assay must also be negative to exclude a monoclonal gammopathy
  \item Immunoparesis without visible M-protein: Suppression of normal immunoglobulins (hypogammaglobulinemia) with a negative SPEP but positive IFE or abnormal FLC ratio may indicate oligosecretory myeloma. The FLC assay is critical in this setting
  \item Hypogammaglobulinemia (decreased gamma globulin fraction <0.7 g/dL): Indicates humoral immunodeficiency (CVID, secondary immunodeficiency from drugs or malignancy, nephrotic syndrome, protein-losing enteropathy) or immunoparesis from an underlying plasma cell dyscrasia
  \item Normal SPEP does not exclude AL amyloidosis, as the monoclonal protein is often very small (<0.5 g/dL and invisible on SPEP); IFE and FLC are more sensitive
\end{itemize}

\subsection*{Positive SPEP (M-Protein Detected)}
\begin{itemize}
  \item M-protein <1.5 g/dL (IgG or IgA) on SPEP with IFE confirmation: Most consistent with monoclonal gammopathy of undetermined significance (MGUS). MGUS is found in 3--5\% of adults over age 50 and up to 7--8\% over age 70. Risk of progression to myeloma is approximately 1\% per year. Risk factors for progression: M-protein $\geq$1.5 g/dL, non-IgG isotype (IgA or IgM), and abnormal FLC ratio \cite{kyle2010monoclonal}
  \item M-protein 1.5--3.0 g/dL: May represent MGUS (with low bone marrow plasma cells and no CRAB criteria) or smoldering multiple myeloma (with $\geq$10\% bone marrow plasma cells). Further workup is determined by the bone marrow biopsy results, FLC assay, and presence or absence of end-organ damage
  \item M-protein $\geq$3.0 g/dL: Multiple myeloma if accompanied by $\geq$10\% bone marrow clonal plasma cells, or smoldering multiple myeloma if no CRAB criteria or myeloma-defining events are present. In the absence of end-organ damage, the diagnosis of smoldering myeloma requires close monitoring (every 3--6 months) because the risk of progression to symptomatic myeloma is approximately 10\% per year for the first 5 years
  \item M-protein >5.0 g/dL: Almost always represents multiple myeloma with a significant tumor burden
  \item IgM monoclonal gammopathy on SPEP/IFE: Waldenstrom macroglobulinemia (IgM M-protein with bone marrow infiltration by lymphoplasmacytic lymphoma). IgM MGUS, if the M-protein is low and there is no marrow infiltration. Both IgM MGUS and Waldenstrom macroglobulinemia carry a risk of hyperviscosity syndrome when the IgM level exceeds 4--5 g/dL due to the large pentameric structure of IgM
  \item Double monoclonal gammopathy (biclonal gammopathy): Two distinct M-proteins on SPEP (e.g., IgG kappa and IgA lambda). Occurs in approximately 2--5\% of monoclonal gammopathies and represents two separate plasma cell clones
  \item Polyclonal hypergammaglobulinemia (broad, diffuse increase in gamma globulins on SPEP, NOT a discrete M-spike): Seen in chronic liver disease (autoimmune hepatitis, primary biliary cholangitis, alcoholic cirrhosis), chronic infections (HIV/AIDS, chronic hepatitis B and C, tuberculosis), autoimmune diseases (SLE, Sjogren syndrome, sarcoidosis, rheumatoid arthritis), and Castleman disease. Immunofixation shows a polyclonal pattern with both kappa and lambda staining
  \item Monoclonal free light chain only (negative SPEP but positive IFE for free light chain): Light chain multiple myeloma. The serum free light chain assay detects the clonal free light chains; the SPEP may show no M-protein because intact immunoglobulins are not produced
  \item Changes in SPEP pattern with therapy: A complete response (CR) requires negative SPEP and negative IFE (the M-protein has disappeared). Very good partial response (VGPR) requires a $\geq$90\% reduction in M-protein. Partial response (PR) requires $\geq$50\% reduction. Progressive disease requires $\geq$25\% increase from nadir (or an absolute increase $\geq$0.5 g/dL)
\end{itemize}

\section*{Normal Results}
\begin{itemize}
  \item \textbf{Serum protein electrophoresis (SPEP) --- normal pattern:}
  \begin{itemize}
    \item Total protein: 6.4--8.3 g/dL
    \item Albumin: 3.5--5.0 g/dL (52--65\% of total protein)
    \item Alpha-1 globulins: 0.1--0.4 g/dL (2--5\%)
    \item Alpha-2 globulins: 0.4--1.0 g/dL (6--12\%)
    \item Beta globulins: 0.5--1.2 g/dL (7--15\%)
    \item Gamma globulins: 0.7--1.6 g/dL (10--20\%)
    \item No monoclonal band/spike detected
    \item Immunofixation: No monoclonal protein; polyclonal kappa and lambda pattern
  \end{itemize}
  \item \textbf{Immunofixation electrophoresis (IFE):} Polyclonal pattern with normal distribution of kappa and lambda light chains (kappa/lambda ratio approximately 2:1). Absence of monoclonal bands with any heavy chain (gamma, alpha, mu, delta, epsilon) or light chain (kappa, lambda) antiserum
  \item \textbf{MGUS (Monoclonal Gammopathy of Undetermined Significance) classification:} M-protein <3 g/dL on SPEP, <10\% clonal plasma cells in bone marrow, AND absence of CRAB criteria or myeloma-defining events. Risk of progression to multiple myeloma or related disorder: 1\% per year
  \item \textbf{Smoldering multiple myeloma:} M-protein $\geq$3 g/dL (IgG or IgA) or urinary M-protein $\geq$500 mg/24 hours, OR $\geq$10--59\% bone marrow plasma cells, AND absence of CRAB criteria or myeloma-defining events. Risk of progression: 10\% per year for first 5 years
  \item \textbf{Multiple myeloma (diagnostic criteria):} M-protein $\geq$3 g/dL AND $\geq$10\% clonal bone marrow plasma cells, OR clonal bone marrow plasma cells $\geq$60\%, OR biopsy-proven plasmacytoma PLUS one or more myeloma-defining events (CRAB criteria or SLiM criteria)
\end{itemize}

\section*{Follow-up Tests}
\begin{itemize}
  \item \textbf{Serum immunofixation electrophoresis (IFE):} Performed automatically in reflex mode when an abnormality (suspected M-spike) is seen on SPEP, or ordered separately for specific clinical indications. IFE confirms the clonality of a suspected M-protein and identifies the immunoglobulin heavy chain class and light chain type
  \item \textbf{Serum free light chain (FLC) assay:} An essential companion test to SPEP and IFE. The combination of SPEP, IFE, and FLC is the IMWG-recommended initial screening panel for monoclonal gammopathies. The FLC assay detects monoclonal free light chains that may not be visible on SPEP (sensitivity for light chain disease is far higher than SPEP alone)
  \item \textbf{24-hour urine protein electrophoresis (UPEP) and urine immunofixation:} Detects and quantifies Bence Jones protein (monoclonal free light chains) in the urine. Important for quantifying heavy proteinuria in AL amyloidosis and light chain cast nephropathy
  \item \textbf{Bone marrow biopsy and aspiration:} Quantifies the percentage of clonal plasma cells. Flow cytometry for plasma cell clonality (CD38+ CD138+ with cytoplasmic kappa or lambda restriction). FISH for high-risk cytogenetic abnormalities [del(17p), t(4;14), t(14;16), t(14;20), 1q gain/ampl]
  \item \textbf{Quantitative immunoglobulins (IgG, IgA, IgM):} Provides the absolute concentration of each immunoglobulin class. Used to confirm immunoparesis (suppression of uninvolved immunoglobulins) and to monitor intact immunoglobulin myeloma where SPEP quantitation may be inaccurate at very low or very high M-protein concentrations
  \item \textbf{Skeletal survey (whole-body low-dose CT, FDG-PET/CT, or whole-body MRI):} Detection of lytic bone lesions, pathologic fractures, and extramedullary plasmacytomas. Whole-body MRI is the most sensitive modality for diffuse bone marrow infiltration
  \item \textbf{Serum calcium, creatinine, complete blood count:} Assessment for the CRAB criteria: hyperCalcemia, Renal insufficiency, Anemia, Bone lesions
  \item \textbf{Serum viscosity:} If symptomatic hyperviscosity is suspected (headaches, blurred vision, epistaxis, neurologic symptoms) in the setting of a high M-protein, particularly IgM >4--5 g/dL, IgA polymers, or IgG3
  \item \textbf{Serum albumin, beta-2 microglobulin, and LDH:} Components of the Revised International Staging System (R-ISS) for multiple myeloma
\end{itemize}

\section*{References}
\bibliographystyle{unsrtnat}
\bibliography{references}

\clearpage
